%%%%%%%%%%  scan idx 203  |  folio 190  %%%%%%%%%%
ainsi d'%
ailleurs que les inclusions $f_{i_0\ldots i_p}:U_{i_0\ldots i_p}=U_{i_0}\cap\ldots\cap U_{i_p}\longrightarrow S$. Soit $E\in\Ob\Qcoh(S)$, et considérons la résolution de $E$ définie par le recouvrement $(U_i)$ (T.F. II 5.2) :
\begin{equation}\tag{*}
 0\longrightarrow E\longrightarrow \coprod f_{i*}f_i^*(E)\longrightarrow\ldots\longrightarrow
 \coprod f_{i_0\ldots i_p*}f_{i_0\ldots i_p}^*(E)\longrightarrow .
\end{equation}
Elle fournit une suite spectrale
\begin{equation}\tag{**}
 E_1^{pq}=\coprod_{i_0\ldots i_p}R^qQ\bigl(f_{i_0\ldots i_p*}f_{i_0\ldots i_p}^*(E)\bigr)
 \Longrightarrow R^{*}Q(E).
\end{equation}
Comme $f_{i_0\ldots i_p}$ est affine, on a d'après (3.4) :
\begin{equation}\tag{***}
 R^qQ\bigl(f_{i_0\ldots i_p*}f_{i_0\ldots i_p}^*(E)\bigr)
 \simeq f_{i_0\ldots i_p*}R^qQ\bigl(f_{i_0\ldots i_p}^*(E)\bigr)
\end{equation}
pour tout $q$, d'où en particulier
\[
 E_1^{pq}=0\quad\text{pour }q>0,
\]
d'après le cas affine, déjà traité. On déduit de (**) des isomorphismes
\[
 R^nQ(E)\simeq H^n(QE^\bullet)\quad\text{pour }n\geq0,
\]
où $E^\bullet$ est la résolution de $E$ définie par la suite exacte (*). Mais d'après (***) on a $QE'{}^n\xrightarrow{\sim} E'{}^n$ pour $n\geq0$, d'où finalement
\[
 R^nQ(E)=0\quad\text{pour }n>0,
\]
et $Q(E)\simeq E$, cqfd.

\emph{Preuve de} (3.5.2). Par la suite spectrale
\[
 R^pQ(H^q(E))\Longrightarrow R^{*}Q(E),
\]
